% Source: https://artofproblemsolving.com/community/c3h2096247_relationship_between_sides_in_a_405090
% Stable identity: aops-topic-2096247-q01
\textbf{(a)} Take the hypotenuse $AB=1$, with the right angle at $C$, $\angle A=40^\circ$ and $\angle B=50^\circ$.

\begin{asy}
usepackage("amsmath");
size(4.5cm);
defaultpen(fontsize(10pt));
pair C=(0,0), A=(cos(radians(40))*5,0), B=(0,sin(radians(40))*5);
draw(A--B--C--cycle,linewidth(0.9));
draw((0.3,0)--(0.3,0.3)--(0,0.3));
label("$C$",C,SW); label("$A$",A,SE); label("$B$",B,N);
label("$40^\circ$",A+(-0.85,0.25),fontsize(8pt)); label("$50^\circ$",B+(0.42,-0.95),fontsize(8pt));
label("$1$",(A+B)/2,NE); label("$\sin40^\circ$",(B+C)/2,W); label("$\cos40^\circ$",(A+C)/2,S);
label("\mbox{}",(min(currentpicture,user=true)+max(currentpicture,user=true))/2,black); // keeps the text after the figure black when LaTeX typesets the labels inline
\end{asy}

The leg opposite $A$ is $BC=\sin40^\circ$ and the leg opposite $B$ is $AC=\cos40^\circ=\sin50^\circ$. So
\[\text{shorter leg}:\text{longer leg}:\text{hypotenuse}=\sin40^\circ:\sin50^\circ:1\approx0.643:0.766:1\approx1:1.192:1.556 .\]

\textbf{(b)} Using $\cos2\theta=2\cos^2\theta-1$ and $\sin2\theta=2\sin\theta\cos\theta$,
\[\cos3\theta=\cos2\theta\cos\theta-\sin2\theta\sin\theta=\left(2\cos^2\theta-1\right)\cos\theta-2\left(1-\cos^2\theta\right)\cos\theta=4\cos^3\theta-3\cos\theta .\]
With $\theta=40^\circ$, $\cos120^\circ=-\dfrac{1}{2}$, so $4c^3-3c=-\dfrac{1}{2}$, that is, $\boxed{8c^3-6c+1=0}$.

\emph{Remark.} So, unlike $\sin30^\circ=\dfrac{1}{2}$ and $\cos30^\circ=\dfrac{\sqrt3}{2}$, the side ratio of a $40$-$50$-$90$ triangle involves a root of a cubic equation. One can show that $\cos40^\circ$ cannot be written using only rational numbers and square roots (this is why a regular $9$-gon cannot be constructed with ruler and compass), so there is no ratio like $1:\sqrt3:2$ for this triangle.
